\section{Introduction}\label{sec:intro}

\subsection{Context and scope}\label{sec:context}

Let $d \ge 2$, let $b = \clog{d}$ be the canonical entry bit-width, and let
$L = d^2 b = \Theta(d^2\log d)$ be the bit-length of a $d \times d$ matrix of $b$-bit entries laid
out row by row on the read-only input tape of a deterministic multitape Turing machine with $k$
read-write work tapes and a one-way write-only output tape
(\cite[Sections~2.1--2.2 and Definition~2.1]{FirstPaper2026}).
Statement~(T) (\cite[Definition~2.1]{FirstPaper2026}, \leanname{Transpose.LinearTimeTransposable})
asserts that some such machine transposes all $d \times d$ matrices of $b$-bit entries within
$O(L)$ steps; Conjecture~9.1 of \cite{FirstPaper2026} (the Tape Coding Conjecture) conjectures, in a
robust co-entropy form, that no causal code solves complete exchange at the profile capacities on the
block networks of linear-length tape diagrams, and its refutation of Statement~(T) is Lean-verified
there (\cite[Theorem~9.2]{FirstPaper2026}). By Harvey and
van der Hoeven \cite[Corollary~6.9]{HvdH2025}, if Statement~(T) is false then no linear-time
multitape integer multiplier exists either.

The synthesis paper \emph{Order, Symmetry and the Swap} \cite{SeventeenthPaper2026} organized the
algebraic obstacles studied across the series into two one-stage mechanisms on $\Mat_{d\times d}(F)$:
an order bound for two-term operators $A_1 X B_1^{\T} + A_2 X B_2^{\T} = X^{\T}$ giving dimension at
most $d(d+1)/2$ (\cite[Theorem~5.1 and Corollary~5.5]{SeventeenthPaper2026}; coordinate form of
Theorem~5.1 Lean-verified, Corollary~5.5 proved in text),
and a translate-pigeonhole bound (Theorem~TP, \cite[Definition~6.3 and Theorem~6.4]{SeventeenthPaper2026},
Lean-verified)
for class-uniform rules in which column $c$ of the output is computed by applying one of $M$ shared
linear decoders to the column-rotated matrix $\rot_c(X)_{i,j} = X_{i, j+c}$. As emphasized in
\cite[Sections~4.2 and~4.6]{SeventeenthPaper2026}, abstract storage-code models without shared rules
admit high-rate network coding, whereas the class-uniform model captures translation covariance along
a single regular sweep.

Multi-stage compositions $K_t \cdots K_1$ of class-uniform hops were first examined for interval
classes in \cite[Remark~6.14 and Section~7.4]{SeventeenthPaper2026} (where $M$ interval classes per
hop gave $\dim_F\ker(\mathrm{sw} - K_t \cdots K_1) \le (d^2 + d + tMd)/2$). In the sequel
\emph{Capacity and Composition of Shared Transposition Decoders}
\cite[Section~6]{SharedDecoders2026} (where all results cited in this paper are Lean-verified),
Brende and GPT-6 Astra studied exact and approximate \emph{two-stage
factorizations} $K_2 \circ K_1$ of linear class-uniform operators on $\Mat_{d\times d}(F^b)$ for
arbitrary column partitions, where
$K_1$ has $M_1$ column classes and $K_2$ has $M_2$ column classes:
\begin{itemize}[leftmargin=1.5em]
\item For exact factorizations $K_2 \circ K_1 = \tau$ over any field $F$,
  Theorem~6.3 of \cite{SharedDecoders2026} reduced two-stage composition to a common-message subspace
  inside a cyclic channel module and extracted a monochromatic arithmetic progression in the stage-1
  partition, proving $d \le M_2\bigl(\mathsf{W}(M_1, M_2 + 1) - 1\bigr)$, where $\mathsf{W}$ denotes
  a van der Waerden threshold; at $M_1, M_2 \le 2$, Corollary~6.4 of \cite{SharedDecoders2026} used
  $\mathsf{W}(2,3) \le 9$ to obtain the small-parameter bound $d \le 16$.
\item For approximate factorizations over any field $F$ with $(d : F) \ne 0$,
  Theorem~6.6 of \cite{SharedDecoders2026}
  proved that for all $d \ge N_0$, $r := \rank_F(\tau - K_2 \circ K_1)$ satisfies
  \[
    r \;\ge\; \frac{b k (d - N_0 + 1)}{2},
    \qquad
    \ell = (2M_2 + 1)k - 1,\quad N_0 = \mathsf{W}(M_1, \ell),
  \]
  showing that $r / (bd) \to \infty$ as $d \to \infty$ for every fixed pair $(M_1, M_2)$, using
  qualitative finiteness of van der Waerden numbers via the Hales--Jewett theorem
  \cite{HalesJewett1963} without substituting a quantitative upper bound for $\mathsf{W}(M_1, \ell)$.
\end{itemize}

This paper is primarily a sequel to \cite{SharedDecoders2026} and \cite{SeventeenthPaper2026} in
that same linear algebraic model. We state the model boundary up front: \emph{no multitape Turing
machine is known to decompose into one- or two-stage linear class-uniform stages}, and none of the
results of this paper decides Statement~(T) or bridges general adaptive or oblivious machines to
linear shared decoders. Throughout this paper, ``Ramsey-free'' means avoiding van der Waerden and
Hales--Jewett progression numbers by substituting Cauchy--Schwarz density cube extraction
\cite{Szemeredi1969,GRS1990} and K\H{o}v\'ari--S\'os--Tur\'an extremal graph bounds
\cite{KovariSosTuran1954}. Within the linear class-uniform and swap-code models, exploiting
\emph{linearity of the shared decoders} makes progress on three open questions from prior work: in
two-stage composition (\cite[Section~8]{SharedDecoders2026}), replacing arithmetic progressions (which
translate-covering $\Z_d$ by a Behrend set \cite{Behrend1946} limits to length $< 3$ at
$M_1 = 2^{O(\sqrt{\log d})}$, and whose non-wrapping monochromatic length in $\{0, \dots, d-1\}$ is
$O(\log d)$ already under a random $2$-coloring \cite{GRS1990,AlonSpencer2016}) by \emph{affine Hilbert
cubes} (over every field) and \emph{shifted monochromatic sumsets} (when $\operatorname{char}(F) \nmid d$,
i.e.\ $(d : F) \ne 0$) eliminates five-fold exponential van der Waerden dependencies \cite{Gowers2001} for
$M_1 \ge 2$; in one-stage linear capacity (\cite[Problem~9.1(a)]{SeventeenthPaper2026};
\cite[Section~3 and Appendix~C]{SharedDecoders2026}), when $\operatorname{char}(F) \nmid d$,
run-extension equations at the right endpoints of monochromatic cyclic runs push strictly below the
$\frac{41}{64}$ bound of \cite[Theorem~3.4]{SharedDecoders2026} and the $\frac{5}{8} - o(1)$ floor
for the nested-component certificates of \cite[Theorem~C.6]{SharedDecoders2026} (whose isolation
variant \cite[Theorem~C.7]{SharedDecoders2026} has floor $\frac{9}{16} - o(1)$); and on \cite[Problem~9.1(b)]{SeventeenthPaper2026} (where run-extension alone gives
$\le 26$), a coupled six-mode Fourier analysis determines $D_1(\{0,1,3\},\{2,4,5\}) = 25$ on $\Z_6$
outside characteristics~$2$ and~$3$.

\subsection{Summary of results}\label{sec:results-overview}

\begin{enumerate}[label=\arabic*.]
\item \textbf{Ramsey-free cutoff for exact two-stage linear composition (\Cref{sec:cubes}).}
  Over an \emph{arbitrary} field $F$, \Cref{lem:affine-cube} generalizes
  \cite[Lemma~6.1]{SharedDecoders2026} from powers of a single endomorphism to ordered products
  $R_{k-1}^{\eps_{k-1}} \cdots R_0^{\eps_0}$ of $k$ endomorphisms $R_0, \dots, R_{k-1} \in \End_F(A)$:
  if $W \le A$ admits an injective linear readout invariant under those $2^k$ products, then
  \[
    (k + 1)\dim_F W \;\le\; \dim_F A \;+\; \sum_{t=0}^{k-1} \dim_F\ker(R_t - \mathrm{id}_A).
  \]
  Extracting a monochromatic affine $M_2$-cube with small steps $1 \le s_t \le \lfloor d/(2M_2^2)\rfloor$
  by the classical Cauchy--Schwarz density cube lemma (\Cref{lem:cs-cube};
  \cite{Szemeredi1969,GRS1990}) yields \Cref{thm:2a}: whenever $b, M_1, M_2 \ge 1$ and two $F$-linear
  class-uniform stages $K_1, K_2$ on $\Mat_{d\times d}(F^b)$ with $M_1$ and $M_2$ classes satisfy
  $K_2 \circ K_1 = \tau$, one has
  \[
    d \;\le\; \max\!\Bigl(8 M_1 M_2^2,\; (2 M_1)^{2^{M_2}}\Bigr).
  \]
  For fixed $M_2 \ge 2$, this cutoff is polynomial in $M_1$---ruling out exact factorizations whenever
  $M_1(d) = o\bigl(d^{1/2^{M_2}}\bigr)$, including partitions of $\Z_d$ into
  $M_1(d) = 2^{O(\sqrt{\log d})}$ $3$-AP-free classes obtained by translate-covering $\Z_d$ with a
  Behrend set \cite{Behrend1946}---whereas every progression-based cutoff from
  \cite[Theorem~6.3]{SharedDecoders2026} is at least $M_2\bigl(\mathsf{W}(M_1, 3) - 1\bigr)$, which is
  superpolynomial in $M_1$. At $M_2 = 1$, $M_1 = 1$, and small constant class counts, however,
  \cite[Theorem~6.3 and Corollary~6.4]{SharedDecoders2026} are stronger: they give $d \le M_1$ at
  $M_2 = 1$, $d \le M_2^2$ at $M_1 = 1$, and $d \le 16$ at $M_1 = M_2 = 2$, where \Cref{thm:2a} gives
  $\max(8M_1, 4M_1^2)$, $\max(8M_2^2, 2^{2^{M_2}})$, and $256$. Two finite controls accompany
  \Cref{thm:2a}: at $d = 13, q = 2$ on the $3$-AP-free $2$-cube $J_{13} = \{0,1,4,5\} \subseteq
  \Z_{13}$, \Cref{ex:d13-witness} shows $9 \le \max\dim_F W \le 10$ over fields with
  $(13 : F) \ne 0$ (refuting a naive $2^k$-multiplier bound $\dim_F W \le 8$); and at $d = 4 = M_1 M_2$
  on $p_1 = p_2 = [0,0,1,1]$, \Cref{ex:d4-coded} constructs exact linear coded factorizations
  $K_2 \circ K_1 = \tau$ for all lane widths $b \ge 0$ over every field with $(5 : F) \ne 0$ (leaving
  characteristic~$5$ open for $b \ge 1$), while showing that entry-copy and bit-lane-copy
  factorizations fail for $b \ge 1$.

\item \textbf{Log-power error-rank lower bound for approximate two-stage composition (\Cref{sec:sumsets}).}
  Over any field $F$ with $(d : F) \ne 0$, we pass to a splitting field carrying an injective
  additive character $\chi : \Z_d \to E^\times$ and replace progressions of length $(2M_2 + 1)k - 1$
  by $(2M_2)$-fold monochromatic sumsets $\widetilde{J}_0 + \widetilde{B}_0 + \cdots +
  \widetilde{B}_{2M_2 - 2}$ with $\abs{\widetilde{J}_0} = k$, $0 \in \widetilde{B}_t$, and
  $\abs{\widetilde{B}_t} = k + 1$. Because $\bigl(\chi(\nu c)\bigr)_{c \in \widetilde{J}_0, 0 \le \nu < k}$
  is a Vandermonde matrix in the distinct scalars $\chi(c)$ ($c \in \widetilde{J}_0$), stacked
  evaluation over any $k$-element set $\widetilde{J}_0$ is injective on the frequency-$< k$ band
  (\Cref{lem:vandermonde-subset}), while each $(k+1)$-element summand $\widetilde{B}_t$ supports a
  nonzero polynomial $P_t$ of degree at most $\diam_{\N}(\widetilde{B}_t)$ annihilating the band
  (\Cref{lem:sumset-poly}). Combining the resulting rank inequality
  $2\rank_F(\tau - K_2 \circ K_1) \ge bk\bigl(d - \sum_t \diam_{\N}(\widetilde{B}_t)\bigr)$
  (\Cref{thm:sumset-rank}) with iterated K\H{o}v\'ari--S\'os--Tur\'an sumset extraction
  (\Cref{lem:kst-sumset}, with $k = O(\log d)$) yields \Cref{thm:2b}: for all $b, M_1, M_2 \ge 1$ and
  all $d \ge D_0(M_1, M_2)$ with $(d : F) \ne 0$,
  \[
    \rank_F(\tau - K_2 \circ K_1)
    \;\ge\;
    c(M_1, M_2)\,b\,d\,(\log d)^{\frac{1}{2M_2 - 1}},
  \]
  and in particular $\Omega(b\,d\log d)$ when $M_2 = 1$ (\Cref{cor:2b-m2-one}). For $M_1 \ge 2$ and
  sufficiently large $d$, this replaces the iterated-logarithmic growth rates obtained from
  \cite[Theorem~6.6]{SharedDecoders2026} via Gowers' van der Waerden bound \cite{Gowers2001} by an
  explicit log-power rate. By contrast, at $M_1 = 1$, $\mathsf{W}(1, \ell) = \ell$ already gives
  $\Omega(bd^2/M_2)$ for all $d \ge 4M_2 + 1$ in \cite[Theorem~6.6]{SharedDecoders2026}; at
  $M_1 = 2, M_2 = 1$, $k = 2$ with $\mathsf{W}(2, 5) = 178$ \cite{StevensShantaram1978} gives
  $b(d - 177)$, which \Cref{thm:2b} ($c(2,1) = \frac{1}{144\log 32}$) beats only for
  $d \gtrsim 10^{216}$ (and the formalized threshold at $(2,1)$ satisfies
  $D_0(2,1) = 32^{36} + 2 \approx 1.53 \times 10^{54}$). \Cref{ex:d12-band} constructs an explicit
  $24$-dimensional shift-covariant $4$-channel band witness ($q = 4$; formalized at $b = 2, M_2 = 1$
  where $bkd = 48$) over $\C$ at $d = 12, k = 2$: under the text-only parameter reading
  $(b, M_2) = (1, 2)$ ($24 = bkd$), a single half-density band without any annihilating-polynomial
  step admits $bkd$-dimensional solutions ($r \ge 0$), showing in that reading that the
  annihilating-polynomial steps in \Cref{thm:sumset-rank} cannot all be dropped.

\item \textbf{Run-extension bounds for one-stage linear shared decoders (\Cref{sec:linear-capacity}).}
  In \cite[Theorem~3.4]{SharedDecoders2026}, iterating
  $C \cap \rot_s^{-1}(C) \le H_{\{0,s\}}$ gave $\frac{41}{64}\,bd^2 + \frac{5}{3}\,bd^{3/2}$ for two
  classes, while the nested-component certificates of \cite{SharedDecoders2026} cannot go below
  $(\frac{5}{8} - o(1))bd^2$ on almost every two-class partition (\cite[Theorem~C.6]{SharedDecoders2026};
  their isolation variant \cite[Theorem~C.7]{SharedDecoders2026} has floor $\frac{9}{16} - o(1)$). For
  $F$-linear decoders (extending the $b = 1$ model of \cite[Definition~6.3 and
  Problem~9.1(a)]{SeventeenthPaper2026} to $b$ lanes), evaluating $\rot_s(X) \in \mathcal{X}$ at the
  right endpoint of each maximal monochromatic run along each non-monochromatic cycle of
  $c \mapsto c + s$ supplies an extra linear equation per run. Filtering by $\rot_s$-eigenmodes over
  $\overline{F}$ proves the single-shift bound $B_{\mathrm{run}}(b, p, s)$
  (\Cref{thm:1a-single-shift}) whenever $\operatorname{char}(F) \nmid (d/\gcd(d,s))$, whose value
  $\frac{1}{4}d^2 + \frac{3}{2}d$ (versus $\dim_F H_{\{0,1\}} = d^2$) on the alternating two-coloring
  for even $d = 2m \ge 2$ at $b = 1, s = 1$ is realized over \emph{every} field
  (\Cref{thm:1a-sharpness}, attaining \Cref{thm:1a-single-shift} whenever
  $\operatorname{char}(F) \nmid 2m$). Averaging over shifts proves
  $D_b(p) \le \frac{19}{32}\,bd^2 + \frac{3}{2}\,bd^{3/2}$ for two classes whenever
  $\operatorname{char}(F) \nmid d$ (\Cref{thm:1b-two-class}, non-vacuous and improving on
  \cite[Theorem~3.4]{SharedDecoders2026} for $d \ge 14$, below $\frac{3}{4}\,bd^2$ for $d \ge 93$, and
  with leading coefficient $\frac{19}{32} < \frac{5}{8} < \frac{41}{64}$; positive characteristics
  dividing $d$ and the gap from $\frac{19}{32}$ to $\frac{1}{2}$ remain open) and
  $D_b(p) \le c_M\,bd^2 + \bigl(\frac{5}{4} - \frac{1}{4M}\bigr)bd$ with
  $c_M = 1 - \frac{5}{4M} + \frac{1}{M^2} - \frac{1}{4M^3}$ for $M \ge 2$ at primes
  $d \ne \operatorname{char}(F)$ (\Cref{thm:1b-multi-class}, giving $\frac{19}{32}\,bd^2 + \frac{9}{8}\,bd$
  at $M = 2$). Its leading coefficient $c_M$ beats the flag bound $1 - \frac{1}{2M}$ of
  \cite[Theorem~5.1, Corollary~5.9(b), Remark~5.10, and Problem~9.1(a)]{SeventeenthPaper2026}
  (coordinate form of Theorem~5.1 Lean-verified; Corollary~5.9(b) and Remark~5.10 proved in text)
  for every $M \ge 2$ (whereas the difference-cover bound \cite[Theorem~3.5, eq.~(6)]{SharedDecoders2026}
  beats the flag bound for $M \ge 7$ and its square-root specialization eq.~(7) does so for $M \ge 11$).
  Against the optimized difference-cover bound \cite[Theorem~3.5, eq.~(6)]{SharedDecoders2026}, $c_M$
  has a strictly smaller leading coefficient if and only if $2 \le M \le 40$---improving on
  \cite[Theorem~3.5]{SharedDecoders2026} at primes $d \ne \operatorname{char}(F)$ past finite
  crossovers (from $d \ge 7$ at $M = 2$ up to $d \ge 39{,}163$ at $M = 40$) and, via
  \Cref{thm:1b-two-class}, for composite $d \ge 35$ with $\operatorname{char}(F) \nmid d$ at $M = 2$.
  Conversely, \cite[Theorem~3.5, eq.~(6)]{SharedDecoders2026} has a $1 - \Theta(M^{-1/2})$ leading
  coefficient that is strictly smaller for every $M \ge 41$, and applies to arbitrary composite $d$
  and arbitrary field characteristic.

\item \textbf{Exact dimension $25$ for $\{0,1,3\} \mid \{2,4,5\}$ on $\Z_6$ outside characteristics $2$ and $3$ (\Cref{sec:aperiodic-six}).}
  On \cite[Problem~9.1(b)]{SeventeenthPaper2026}, combining six-mode column-Fourier decomposition on
  $H_{\{0,1\}}$, an antipodal quotient bound $3\dim Q + S_K \le 24$, and kernel-disjointness /
  $4 \times 4$ row elimination (cleared determinant $2^4 \cdot 3^3 \ne 0$) at $S_K \in \{11, 12\}$ proves
  $\dim_F(\mathcal{X} \cap \rot_1^{-1}(\mathcal{X})) \le 14$ (\Cref{lem:aperiodic-six-shift-14}) and
  hence $D_1\bigl(\{0,1,3\},\{2,4,5\}\bigr) = 25$ over every field with
  $\operatorname{char}(F) \notin \{2, 3\}$ (\Cref{thm:aperiodic-six-25}, paired with the lower bound
  of \cite[Theorem~3.2]{ClassUniformMemo2026}); characteristics~$2$ ($25 \le D_1 \le 28$,
  \cite[Theorem~6.4 and Proposition~8.2]{SeventeenthPaper2026}, \cite[Theorem~3.2]{ClassUniformMemo2026})
  and~$3$ ($21 \le D_1 \le 28$ from symmetric matrices and \cite[Theorem~6.4]{SeventeenthPaper2026};
  whether $25$ is attainable in characteristic~$3$ is \cite[Problem~6.1]{ClassUniformMemo2026}) remain
  open.

\item \textbf{Three-term swap codes and one- and two-sided interval-mask bounds (\Cref{sec:swapcodes}).}
  In \cite[Remark~5.6]{SeventeenthPaper2026}, examples of three-term swap codes
  $\{X \in \Mat_{d\times d}(F) : \sum_{t=0}^2 A_t X B_t^{\T} = X^{\T}\}$ exceeding $d(d+1)/2$ were
  given at $d = 3$ ($8 > 6$, over every field) and $d = 4$ ($11 > 10$, checked numerically over
  $\F_2, \F_3, \F_5, \F_7$), with no example known for $d \ge 5$. In
  \Cref{thm:three-term-kronecker,thm:three-term-unified}, we construct a $3 \times 3$ block family
  on $d = 3m$ achieving exact dimension $4m(m+1)$ over every field---exceeding $d(d+1)/2$ at
  $d \in \{3, 6, 9, 12\}$---and pad it to obtain $29 > 28$ at $d = 7$ and $11 > 10$ at $d = 4$ over
  every field; $d = 5$ and every $d \ge 8$ outside $\{9, 12\}$ remain open. On
  \cite[Problem~9.3]{SeventeenthPaper2026}, after recording the one-sided endpoint-refinement baseline
  $(d^2 + d\min(\beta_A, \beta_B))/2$ (\Cref{prop:multiscale-interval}), we prove a
  characteristic-free two-sided bound for two-scale (coarse--fine) families
  (\Cref{thm:two-sided-coarse-fine,cor:two-sided-win}) whose leading coefficient on contiguous fine
  singletons plus one coarse complement,
  $\frac{1}{2} + \frac{1}{4}\min\bigl(\beta(1+\alpha),\alpha(1+\beta)\bigr)$, is strictly below
  one-sided Theorem~TP's $\frac{1}{2} + \frac{1}{2}\min(\alpha,\beta)$
  (\cite[Theorem~6.4]{SeventeenthPaper2026}) for all $\alpha, \beta \in (0,1)$ (giving
  $\frac{11}{16}d^2 + d < \frac{3}{4}d^2$ at $\alpha = \beta = 1/2$ for all even $d \ge 18$), while
  \Cref{prop:walsh-dichotomy} contrasts interval boundaries ($\le 2$) with periodic Walsh digit masks
  ($2^{j+1}$ boundaries, profile $2$ per scale, and operator span dimension $\le (m+1)2^m$).

\item \textbf{Restricted shuffle and row/column-uniform operator pipelines and anti-diagonal untwisting (\Cref{sec:duality}).}
  Across the row/column tensor cut $F_{\mathrm{row}}^d \otimes_F (F_{\mathrm{col}}^d \otimes_F F^b)$,
  Operator Schmidt Rank satisfies $\OSR(\tau) = d^2$ and $\OSR(S) = \OSR(S^{-1}) = 4$ for Stone's
  perfect shuffle and unshuffle \cite{Stone1971,Nielsen2003,Tyson2003} on $d = 2^m \ge 2$
  (\Cref{thm:osr-swap}), forcing $t \ge m = \log_2 d$ shuffle/unshuffle passes (attained by
  $S^m = \tau$) in any pipeline composed solely of $S, S^{-1}$, row-uniform operators $I_d \otimes B$,
  and column-uniform operators $A \otimes I_{db}$ (\Cref{thm:osr-pass}). As we emphasize at the start
  of \Cref{sec:duality}, this applies only to that restricted operator pipeline and \emph{does not}
  answer \cite[Problem~9.5]{SeventeenthPaper2026}: within-row cyclic shear
  $X_{r,c} \mapsto X_{r, (c+r)\bmod d}$ has $\OSR = d$ (\Cref{prop:shear-control}; see also
  \cite[Proposition~6.25(a)]{SeventeenthPaper2026} for affine row-coset stages) yet runs in $O(L)$
  time on a $2$-work-tape Turing machine (\Cref{prop:stone-tm}, proved in the shared library
  \texttt{lean-machine} alongside $3$-tape shuffle and $2$-tape unshuffle). \Cref{sec:antidiag}
  repackages the shift-covariance subspace of \cite[Lemma~3.1]{SharedDecoders2026} (together with the
  compatibility equivalence \leanname{SharedDecoders.isSharedCode_iff_compatible} in its Lean library)
  via an
  explicit \leanname{TensorProduct} isomorphism intertwining $\rot_t$ with $\mathrm{id} \otimes \sigma^t$
  (where $(\sigma^t v)_a = v_{a+t}$, \Cref{thm:anti-diagonal}) and records the two-subspace Grassmann
  relaxation (\Cref{prop:mask-floors}).

\item \textbf{Cut-cone evaluation, random-order set capacity, and negative brackets for two column-class cut reductions (\Cref{sec:sweep}).}
  On frozen trace/cut fibres of \leanname{Transpose.TM}, \Cref{lem:tm-cut-cone} extends the causal-cone
  lemmas of \cite[Lemma~4.4 and Corollary~4.5]{FourthPaper2026},
  \cite[Lemmas~3.3(b) and~4.7(a)]{MergeDepth2026}, and
  \cite{AdaptiveTransfer2026} by instantiating the $(M.k+1)^D$ cone-size bound on the native step DAG
  and constructing depth-$D$ symbolic \lean{CausalConeTree}s directly from
  $M.\delta$ whose syntactic input-cell support has size at most $b(M.k+1)^D$.
  \Cref{thm:random-order-capacity} formalizes in Lean the column-class analogue of the
  $1/(\Delta+1)$ random-order acyclic-subgraph bound of
  \cite[Proposition~4.14(a)--(b)]{FourthPaper2026} (after \cite[Theorem~3]{BBJK2011}). Finally,
  \Cref{rem:sweep-negative} and \Cref{prop:negative-brackets} illustrate that two naive column-class
  cut reductions---discarding cut-touching outputs, or positing a worst-case full-product class count
  across per-column cut slots---produce strictly negative deficit brackets along the control spine
  $u \to u+1$, reinforcing the observation in the commentary on
  \cite[Problem~9.4(b)]{SeventeenthPaper2026} regarding per-position fixed-alphabet recovery.
\end{enumerate}

\subsection{Formalization conventions}\label{sec:conventions}

Every numbered theorem, lemma, proposition, corollary, and example in this paper has a kernel-verified
Lean~4 counterpart (version \texttt{v4.34.1} with Mathlib \cite{Mathlib2026}) across the
\texttt{BoundedArityDecoders} package and the shared \texttt{LeanMachine.ShuffleShear} modules in
\texttt{lean-machine}, depending solely on Lean's three standard axioms (\texttt{propext},
\texttt{Classical.choice}, and \texttt{Quot.sound}), with zero \texttt{sorry}, \texttt{admit},
\texttt{native\_decide}, or custom axioms. Each numbered result carries a \textsf{Lean:} status line
citing proved \texttt{theorem} declarations, followed by an outline of the proof; where a companion
asymptotic formulation, external-theorem specialization, or parameter re-reading is proved in the text
only, the status line marks the result \emph{partially formalized} and states the text-only gap
explicitly (summarized in \Cref{sec:formalization}).
